\documentclass[12pt]{article}
\usepackage{amsmath, amssymb}
\usepackage[margin=1in]{geometry}
\setlength{\parskip}{6pt}
\title{QUBO Formulation: Planar Robot Inverse Kinematics Candidate Selection}
\date{}
\begin{document}
\maketitle

\subsection*{Planar Robot Inverse Kinematics Candidate Selection}

\paragraph{Problem Statement.}
Given a set of $N$ discretized joint-angle candidates for a planar robot, the task is to select exactly one candidate such that the sum of the inverse kinematics (IK) error and a regularization cost associated with the selected candidate is minimized.

\paragraph{Decision Variables.}
Let $N$ be the number of discretized joint-angle candidates. We introduce binary decision variables $x_i \in \{0, 1\}$ for each candidate index $i \in \{0, 1, \dots, N-1\}$. The variable $x_i$ indicates the selection of candidate $i$, where $x_i = 1$ if candidate $i$ is chosen and $x_i = 0$ otherwise.

\paragraph{QUBO Derivation.}
\textbf{Step 1: Objective Function.}
Let $c_i$ denote the combined IK error and regularization cost for candidate $i$. The objective is to minimize the total cost:
\begin{align}
    \min \sum_{i=0}^{N-1} c_i x_i.
\end{align}
Since the QUBO model minimizes the Hamiltonian $H(x)$, we identify the objective contribution as:
\begin{align}
    H_{\text{obj}}(x) = \sum_{i=0}^{N-1} c_i x_i.
\end{align}

\textbf{Step 2: Constraints.}
The problem requires that exactly one candidate be selected. This is enforced by the equality constraint:
\begin{align}
    \sum_{i=0}^{N-1} x_i = 1.
\end{align}

\textbf{Step 3: Penalty Terms.}
We convert the constraint into a penalty term using a positive weight $\lambda$. The penalty Hamiltonian is:
\begin{align}
    H_{\text{pen}}(x) = \lambda \left( \sum_{i=0}^{N-1} x_i - 1 \right)^2.
\end{align}
Expanding the square yields:
\begin{align}
    H_{\text{pen}}(x) = \lambda \left( \sum_{i=0}^{N-1} x_i^2 - 2 \sum_{0 \le i < j \le N-1} x_i x_j + 1 \right).
\end{align}
Using the idempotent property of binary variables, $x_i^2 = x_i$, we simplify the quadratic sum of variables:
\begin{align}
    \sum_{i=0}^{N-1} x_i^2 = \sum_{i=0}^{N-1} x_i.
\end{align}
Substituting this back into the penalty expression gives:
\begin{align}
    H_{\text{pen}}(x) = \lambda \left( \sum_{i=0}^{N-1} x_i - 2 \sum_{0 \le i < j \le N-1} x_i x_j + 1 \right).
\end{align}
Distributing $\lambda$ and separating terms, we obtain the general penalty contribution:
\begin{align}
    H_{\text{pen}}(x) = -\lambda \sum_{i=0}^{N-1} x_i + 2\lambda \sum_{0 \le i < j \le N-1} x_i x_j + \lambda.
\end{align}

\textbf{Step 4: Combined Hamiltonian.}
The total QUBO Hamiltonian $H(x)$ is the sum of the objective and penalty terms:
\begin{align}
    H(x) = H_{\text{obj}}(x) + H_{\text{pen}}(x) = \sum_{i=0}^{N-1} c_i x_i - \lambda \sum_{i=0}^{N-1} x_i + 2\lambda \sum_{0 \le i < j \le N-1} x_i x_j + \lambda.
\end{align}
Grouping the linear terms, we write:
\begin{align}
    H(x) = \sum_{i=0}^{N-1} (c_i - \lambda) x_i + 2\lambda \sum_{0 \le i < j \le N-1} x_i x_j + \lambda.
\end{align}

\textbf{Step 5: Q Matrix and Offset.}
The Hamiltonian is expressed in the standard QUBO form $H(x) = x^T Q x + c$, where $x = [x_0, \dots, x_{N-1}]^T$. Reading off the coefficients from the general expression above, the entries of the symmetric matrix $Q$ and the scalar offset $c$ are:
\begin{align}
    Q_{i,i} &= c_i - \lambda, \quad \forall i \in \{0, \dots, N-1\}, \\
    Q_{i,j} &= 2\lambda, \quad \forall 0 \le i < j \le N-1, \\
    Q_{j,i} &= 2\lambda, \quad \forall 0 \le i < j \le N-1, \\
    c &= \lambda.
\end{align}

\paragraph{Penalty Weight Justification.}
The penalty weight $\lambda$ must be chosen sufficiently large to ensure that any violation of the constraint results in a higher energy than any feasible solution. The maximum possible objective cost for a feasible solution is $\max_{i} c_i$. If the constraint is violated, the penalty term $\lambda (\sum x_i - 1)^2$ takes a value of at least $\lambda$ (since $\sum x_i$ is an integer, any violation implies $|\sum x_i - 1| \ge 1$). Therefore, choosing $\lambda > \max_{i} c_i$ guarantees that the energy of any infeasible solution strictly exceeds the energy of any feasible solution, ensuring the global minimum corresponds to a feasible candidate.

\paragraph{Correctness Argument.}
Minimizing $H(x)$ over $x \in \{0,1\}^N$ recovers the optimal candidate selection. For any feasible solution satisfying $\sum_{i=0}^{N-1} x_i = 1$, the penalty term vanishes, and $H(x) = \sum_{i=0}^{N-1} c_i x_i = c_k$ where $k$ is the unique index with $x_k=1$. Thus, among feasible solutions, $H(x)$ equals the cost of the selected candidate. For infeasible solutions, the penalty term contributes at least $\lambda$. By the choice $\lambda > \max_i c_i$, the energy of any infeasible solution is strictly greater than the cost of the best feasible candidate. Consequently, the global minimum of $H(x)$ must be feasible and corresponds to the candidate with the minimum cost $c_k$.

\paragraph{Illustrative Example.}
Consider a concrete instance with $N=3$ candidates and costs $c_0 = 10$, $c_1 = 5$, and $c_2 = 8$. We choose $\lambda = 100$, which satisfies $\lambda > \max\{10, 5, 8\}$.
The QUBO Hamiltonian becomes:
\begin{align}
    H(x) &= (10 - 100)x_0 + (5 - 100)x_1 + (8 - 100)x_2 + 200 x_0 x_1 + 200 x_0 x_2 + 200 x_1 x_2 + 100 \nonumber \\
    &= -90 x_0 - 95 x_1 - 92 x_2 + 200 x_0 x_1 + 200 x_0 x_2 + 200 x_1 x_2 + 100.
\end{align}
The optimal solution is $x^* = [0, 1, 0]^T$, corresponding to selecting candidate 1. The minimum energy is $H(x^*) = -95 + 100 = 5$, which matches the cost $c_1$.

\end{document}